\chapter{Technical Debt}
\label{ch:technical-debt}

\section{Definition and method}

Technical debt here means maintenance cost accepted in the current codebase:
missing safety nets, duplication, hard-coded values, and structure that will
resist change. Each item has an identifier, an estimate of the interest it
accrues (how much it will slow future work), the evidence, and a recommended
action. Items are ordered by the value of paying them down.

This is distinct from Chapter~\ref{ch:known-issues}, which describes observable
defects. Some findings appear in both where a defect is also a structural
weakness.

\section{Debt register}

\begin{longtable}{p{1.2cm}p{3.6cm}p{1.7cm}p{7.1cm}}
\caption{Technical debt register}\\
\toprule
\textbf{ID} & \textbf{Item} & \textbf{Interest} & \textbf{Evidence and consequence} \\
\midrule

TD-01 & No automated tests anywhere. & High & Absence of any test file. Every
change to the shared core engine or to a template is unverified except by hand;
regressions reach users. The core engine is used by two clients, so a single
mistake is amplified. \\
TD-02 & No linting or formatting configuration. & Medium & No \texttt{eslint},
\texttt{prettier} or \texttt{editorconfig}. Style is upheld by convention only;
diffs carry incidental noise and code quality is inconsistent. \\
TD-03 & No continuous integration on pull requests. & High & Neither workflow runs
on \texttt{pull\_request}. Broken builds and invalid recipes are discovered late
(on merge or at release). See Chapter~\ref{ch:cicd}. \\
TD-04 & Catalog inlined into every client. & Medium & The catalog is duplicated as
\texttt{bundle.json}, per-client generated TypeScript and code inside each bundle
($\approx$240\,KB each). It grows linearly with the recipe count and inflates
both the npm package and the extension. See Chapter~\ref{ch:scalability}. \\
TD-05 & Templates are recompiled on each render. & Low & \texttt{renderTemplate}
calls \texttt{Handlebars.compile} every time. A compiled-template cache is a
small, safe improvement. \\
TD-06 & Hard-coded CLI version string. & Low & Two authorities for the version
number will diverge on release. Read from \texttt{package.json} instead. \\
TD-07 & Hand-written schema validation. & Medium & \texttt{validate.mjs}
re-implements a subset of JSON Schema, so it can drift from
\texttt{schema.json} and it is not exercised in CI. Use a real validator. \\
TD-08 & Duplicated generation logic across the two clients. & High & The CLI
(\texttt{commands/add.ts} and utilities) and the extension
(\texttt{generator.ts}) each implement detection, option resolution, file
writing, manifest handling and primitive resolution independently. Fixes must be
applied twice, and the two paths can diverge. \\
TD-09 & No root README or contribution guide. & Medium & The repository root has
no \texttt{README.md}. Onboarding a contributor requires reading source. See
Chapter~\ref{ch:contribution-guide}. \\
TD-10 & Generated artefacts tracked in version control. & Low & The built
registry (\texttt{registry/}, 25 files) and the offline bundle
(\texttt{packages/registry/bundle.json}) are committed even though they are
derived from \texttt{packages/registry/data/}. The four local \texttt{.vsix}
files are correctly excluded by \texttt{.gitignore} but remain in the working
tree. Derived files in history invite drift and noisy diffs. \\
TD-11 & Registry package lacks a licence declaration. & Low &
\texttt{packages/registry/package.json} has no \texttt{license} field, while the
other packages do. Ambiguous reuse terms. \\
TD-12 & No error-classification or logging convention. & Medium & Errors are
surfaced with ad-hoc strings; there is no shared logging layer or set of error
types. Diagnostics are hard to improve consistently. \\
TD-13 & Silent tolerance of corrupt state. & Low & An unreadable manifest is
treated as empty. Convenient, but it hides corruption. \\
TD-14 & No dependency-update automation or audit. & Medium & No Dependabot,
Renovate or \texttt{pnpm audit} in CI; caret ranges mean drift without signal. \\
TD-15 & Divergent shell execution across platforms. & Low & Generators shell out
for installation; the approach is not wrapped behind a single abstraction,
complicating testing and cross-platform reasoning. \\
\bottomrule
\end{longtable}

\section{The dominant debt: two clients, one engine, duplicated orchestration}

\texttt{@expo-snips/core} correctly centralises \emph{rendering}. What it does
\emph{not} centralise is everything around rendering: project detection, option
precedence, the shadcn enforcement rules, primitive resolution, file writing,
manifest updates and dependency installation. Both the CLI and the extension
implement these separately.

The evidence:

\begin{itemize}
  \item \texttt{packages/cli/src/utils/detectProject.ts} and
        \texttt{packages/vscode-extension/src/project.ts} both probe
        \texttt{package.json} and \texttt{tsconfig.json}.
  \item \texttt{packages/cli/src/utils/primitives.ts} and the download/primitive
        logic in \texttt{packages/vscode-extension/src/generator.ts} both resolve
        primitive recipes.
  \item \texttt{packages/cli/src/utils/registry.ts} and
        \texttt{packages/vscode-extension/src/registry.ts} both fetch and fall
        back.
  \item \texttt{packages/cli/src/utils/manifest.ts} and the manifest functions in
        \texttt{project.ts} both read and write \texttt{.expo-snips.json}.
\end{itemize}

\textbf{Consequence:} a behavioural change (for example, ``shadcn always implies
NativeWind and Lucide'') must be made in more than one place, and there is
currently no test that would catch a divergence.

\textbf{Recommendation (RECOMMENDED):} promote the orchestration layer into
\texttt{@expo-snips/core} (or a new shared package) so that both clients call one
implementation, leaving only UI/terminal concerns in the clients.

\section{Debt heat map}

\begin{table}[H]
\centering
\caption{Debt by theme and interest}
\begin{tabular}{p{5cm}p{1.7cm}p{7.5cm}}
\toprule
\textbf{Theme} & \textbf{Interest} & \textbf{Items} \\
\midrule
Verification (tests, CI, lint) & High & TD-01, TD-03, TD-02, TD-07, TD-14 \\
Shared architecture & High & TD-08, TD-04 \\
Documentation / onboarding & Medium & TD-09, TD-11 \\
Robustness & Medium & TD-12, TD-13 \\
Hygiene & Low & TD-10, TD-06, TD-05, TD-15 \\
\bottomrule
\end{tabular}
\end{table}

\section{Priority}

\begin{enumerate}
  \item \textbf{Introduce CI with validation and build} (TD-03) --- nearest-term,
        protects everything else.
  \item \textbf{Add golden-file tests for the core engine} (TD-01) --- makes the
        shared engine safe to refactor.
  \item \textbf{Unify the orchestration layer} (TD-08) --- removes the highest
        structural interest.
  \item \textbf{Address the output-correctness issues} in
        Chapter~\ref{ch:known-issues} (KI-01, KI-02, KI-03) --- these are product
        promises, not just hygiene.
  \item \textbf{Pay down the remaining low-interest items} opportunistically.
\end{enumerate}

\section{Summary}

\textbf{Status:} the debt is concentrated and identifiable rather than diffuse.
The core rendering engine is small, well-typed and reasonably factored; the debt
lives in verification (no tests, no PR CI) and in duplicated orchestration across
the two clients. Both are addressable without rearchitecting, and both have high
leverage because the engine is shared.
